% Bản nháp phần phụ lục theo dàn bài mới.
% ===========================================================================
\appendix

\chapter{Kiểm thử tính đúng đắn}
\label{app:tests}

Trước khi chạy lưới tham số, 60 phép thử tự động kiểm tra toàn bộ cài đặt. Bài
toán có sẵn nghiệm đóng và có sẵn biểu thức của gradient lẫn Hessian, nên phần
lớn phép thử đối chiếu thẳng với một mốc giải tích chứ không so với một lần chạy
trước đó.

\begin{enumerate}
  \item \textbf{Gradient giải tích so với sai phân trung tâm.} So sánh
    $[\nabla f(w)]_i$ với
    $\bigl(f(w + \eps e_i) - f(w - \eps e_i)\bigr) / (2\eps)$ tại
    $\eps = 10^{-6}$. Sai số tương đối phải dưới $10^{-6}$.
  \item \textbf{Hessian so với sai phân của gradient.} Cùng ngưỡng.
  \item \textbf{Nghiệm đóng là điểm dừng.} $\norm{\nabla f(\wstar)} < 10^{-12}$.
  \item \textbf{Mọi thuật toán cùng hội tụ về một nghiệm.} Trên bài toán tổng hợp
    $n = 100$, $d = 5$, cả mười cấu hình tất định đều cho
    $\norm{w_{\text{cuối}} - \wstar} < 10^{-6}$.
\end{enumerate}

Ngoài ra có bốn phép thử về mặt phương pháp luận, tức chúng kiểm chính cách đo
chứ không kiểm công thức:

\begin{itemize}
  \item gradient descent phải phân kỳ khi $t > 2/L$, vì phép thử ngưỡng ổn định ở
    mục~\ref{sec:mech-step} sẽ vô nghĩa nếu tiêu chí dừng theo đình trệ nuốt mất
    trường hợp này;
  \item SGD với độ dài bước hằng phải dừng lại cách $\fstar$ một khoảng dương;
  \item momentum của phương pháp tăng tốc phải tính theo độ dài bước thực tế, kiểm
    trên một bài toán có điều kiện xấu với line search bật;
  \item ghi lịch sử dày hay thưa phải cho cùng số lần đánh giá hàm mục tiêu, tức
    chi phí theo dõi không lẫn vào chi phí thuật toán.
\end{itemize}

Ba phép thử cuối cùng chạy trên chính nguồn LaTeX chứ không trên mã Python: chúng
kiểm rằng mọi hình và bảng đều có nhãn và được nhắc tới, rằng không hình nào bị
chèn hai lần, và rằng mọi mục trong thư mục tài liệu tham khảo đều được trích dẫn.
Quy tắc về nhãn từng bị vi phạm ở 6 trên 10 bảng và 18 trên 22 hình khi chỉ kiểm
bằng mắt, nên nó được chuyển thành phép thử.

Các phép thử và chính các thí nghiệm đã bắt được ba lỗi có thật. Thứ nhất là cách
nhận diện cột số có đơn vị, vốn biến cột \texttt{model} thành số và làm khối
one-hot mất đi biến định tính nhiều mức nhất. Thứ hai là momentum không nhất quán
với độ dài bước, đẩy hàm mục tiêu lên $1{,}47 \cdot 10^{4}$
(mục~\ref{sec:mech-momentum}). Thứ ba là thiếu tiêu chí dừng theo đình trệ, khiến
một lần chạy tiêu 98,7\% thời gian vào phần đã hết cải thiện
(mục~\ref{sec:recorder}). Cả ba đều để chương trình chạy trót lọt và trả về một
con số trông hợp lý, nên không phép thử nào đối chiếu với mốc giải tích thì chúng
còn nguyên trong kết quả.

\chapter{Ba thuật toán không đưa vào so sánh chính}
\label{app:cut}

Ngoài năm thuật toán ở chương~\ref{sec:algorithms}, nhóm cài thêm Heavy ball,
Newton-CG và Adam. Cả ba đều qua bộ kiểm thử ở phụ lục~\ref{app:tests} nhưng
không vào bảng~\ref{tab:summary}, và bảng~\ref{tab:cut} nói vì sao.

\begin{table}[tp]
  \centering
  \caption{Ba thuật toán đã cài nhưng không đưa vào so sánh chính, đo trên cùng
    một máy với cùng quy ước như bảng~\ref{tab:summary}.}
  \label{tab:cut}
  \begin{tabular}{lrrl}
    \toprule
    Cấu hình & Vòng lặp đạt $10^{-6}$ & Thời gian (s) & Đối chiếu \\
    \midrule
    Heavy ball, tham số tối ưu   & 55        & \num{1.35}  & GD thường cần 20 vòng \\
    Newton-CG, \texttt{cg\_tol} $= 10^{-6}$ & 1 & \num{0.596} & Newton cần \num{0.133} \\
    Adam, toàn bộ dữ liệu        & 115       & \num{3.50}  & GD cần \num{0.465} \\
    Adam, $B = 64$               & không đạt & --          & dừng ở $3{,}1 \cdot 10^{-2}$ \\
    \bottomrule
  \end{tabular}
\end{table}

Heavy ball~\cite{polyak1964heavyball} chỉ khác phương pháp tăng tốc ở thứ tự giữa
bước ngoại suy và bước lấy gradient, nên nhóm cài nó để đặt cạnh Nesterov. Với bộ tham số tối ưu cho hàm bậc
hai, $t = 4/(\sqrt{L}+\sqrt{\mu})^{2}$ và
$\beta = \bigl((\sqrt{\kappa}-1)/(\sqrt{\kappa}+1)\bigr)^{2}$, nó cần 55 vòng lặp
để đạt $10^{-6}$, tức nhiều hơn gradient descent thường. Kết quả này không mâu
thuẫn với lý thuyết mà là ví dụ thứ ba của cùng một hiện tượng ở
mục~\ref{sec:theory-early}: cả hai tham số đều tính từ $\kappa$ nên chúng chỉnh
cho trường hợp phổ trải đều, còn phổ của bài toán này thì dồn về đầu lớn.

Newton-CG giải hệ Newton bằng gradient liên hợp thay vì phân rã Cholesky, nên mỗi
bước chỉ cần một tích Hessian nhân vector với chi phí $\bigO(nd)$ và không bao giờ
phải lập ma trận $d \times d$. Ở $d = 280$ nó chậm hơn Newton khoảng 10 lần
(mục~\ref{sec:mech-cost}), nhưng nó tồn tại đúng vì trường hợp $d$ lớn, tức đúng
trường hợp mà báo cáo này không chạm tới.

Adam~\cite{kingma2015adam} chia bước theo căn bậc hai của moment bậc hai, cơ chế chỉ phát huy khi thang
gradient lệch nhau nhiều giữa các tọa độ. Bước chuẩn hóa cột ở
chương~\ref{sec:data} đã loại đúng chênh lệch ấy, nên trên bài toán này Adam
không có lợi thế nào: bản toàn bộ dữ liệu cần 115 vòng lặp so với 20 của gradient
descent, còn bản $B = 64$ không đạt $10^{-6}$ trong 400 vòng.

\chapter{Quy đổi hàm mục tiêu của scikit-learn}
\label{app:conversion}

Phép so ở mục~\ref{sec:compare-sklearn} chỉ có nghĩa khi hai bên tối ưu cùng một
hàm, và ba bộ ước lượng dùng ba cách chuẩn hóa hệ số khác nhau.

\texttt{Ridge} cực tiểu hóa $\norm{Xw-y}_{2}^{2} + \alpha \norm{w}_{2}^{2}$, tức
không có hệ số $\tfrac{1}{2n}$ lẫn $\tfrac{\lambda}{2}$ của
\eqref{eq:objective}. Nhân \eqref{eq:objective} với $2n$ được
$\norm{Xw-y}_{2}^{2} + \lambda n \norm{w}_{2}^{2}$, và vì nhân hàm mục tiêu với
một hằng số dương không đổi nghiệm nên
%
\begin{equation}
  \label{eq:alpha-ridge}
  \alpha_{\texttt{Ridge}} = \lambda n .
\end{equation}
%
Với $\lambda = \num{0.01}$ và $n = \num{160000}$ thì $\alpha = \num{1600}$.

\texttt{SGDRegressor} với \texttt{loss="squared\_error"} và \texttt{penalty="l2"}
cực tiểu hóa trung bình mất mát trên các mẫu cộng $\tfrac{\alpha}{2}\norm{w}^{2}$,
tức trùng \eqref{eq:objective} ngay khi $\alpha = \lambda$. \texttt{LinearRegression}
không có số hạng phạt nên nó ứng với $\lambda = 0$ và giải một bài toán khác;
báo cáo đưa nó vào để có mốc tham chiếu chứ không để so tốc độ.

Nhóm kiểm \eqref{eq:alpha-ridge} bằng số chứ không bằng suy luận trên giấy. Lấy
$\hat{w}$ do \texttt{Ridge} với solver \texttt{cholesky} trả về, khoảng cách tới
nghiệm đóng \eqref{eq:closed-form} là $1{,}5 \cdot 10^{-15}$, sai số tương đối
$1{,}9 \cdot 10^{-15}$, và $f(\hat{w}) - \fstar$ bằng đúng 0 trong số học dấu
phẩy động. Phép kiểm này nằm trong bộ kiểm thử chứ không phải một dòng in ra để
đọc bằng mắt, vì mọi con số ở bảng~\ref{tab:sklearn} sẽ vô nghĩa nếu nó thất bại.

\chapter{Dẫn xuất các công thức}
\label{app:derivations}

\section{Từ hàm mục tiêu tới nghiệm đóng}

Viết $r(w) = Xw - y$. Vì
$\norm{r}_{2}^{2} = r\trans r$ và $\partial r / \partial w = X$, quy tắc dây
chuyền cho $\nabla_w \tfrac{1}{2n}\norm{r}^{2} = \tfrac{1}{n} X\trans r$. Đạo hàm
của $\tfrac{\lambda}{2}\norm{w}^{2}$ là $\lambda w$, nên cộng lại được vế thứ
nhất của \eqref{eq:derivatives}. Lấy đạo hàm lần nữa, $X\trans r$ tuyến tính theo
$w$ với ma trận $X\trans X$, nên $\nabla^{2} f = \tfrac{1}{n} X\trans X + \lambda I$
và nó không chứa $w$.

Cho gradient bằng không được
$\bigl(\tfrac{1}{n} X\trans X + \lambda I\bigr) w = \tfrac{1}{n} X\trans y$. Ma
trận vế trái xác định dương với mọi $\lambda > 0$ vì $X\trans X$ nửa xác định
dương, nên hệ có nghiệm duy nhất, đó là \eqref{eq:closed-form}.

\section{Hằng số Lipschitz là chuẩn phổ của Hessian}

Với $H = \nabla^{2} f$ hằng, $\nabla f(u) - \nabla f(v) = H(u-v)$ với mọi $u, v$.
Do đó $\norm{\nabla f(u) - \nabla f(v)} = \norm{H(u-v)} \le \norm{H}_{2} \norm{u-v}$,
tức \eqref{eq:lipschitz} đúng với $L = \norm{H}_{2}$. Chặn này đạt được khi
$u - v$ nằm theo hướng riêng ứng với trị riêng lớn nhất, nên không lấy $L$ nhỏ
hơn được. Vì $H$ đối xứng nửa xác định dương cộng $\lambda I$, chuẩn phổ của nó
bằng trị riêng lớn nhất, và ta được \eqref{eq:constants}.

\section{Hằng số trơn của một lô}

Gọi $L_{\max}$ là hằng số trơn lớn nhất trong các mất mát một mẫu, ở đây
$\norm{x_i}^{2} + \lambda$. Mất mát trên một lô $B$ mẫu lấy ngẫu nhiên là trung
bình của $B$ hàm, và với $B$ mẫu độc lập thì phần phương sai giữa các mẫu chia
cho $B$ còn phần chung giữ nguyên, nên hằng số trơn nội suy giữa hai đầu:
$L_B \to L_{\max}$ khi $B = 1$ và $L_B \to L$ khi $B = n$. Dạng nội suy tuyến
tính theo $1/B$ cho \eqref{eq:minibatch-smoothness}. Đây là xấp xỉ dùng làm quy
tắc chọn bước chứ không phải đẳng thức chặt; điều duy nhất báo cáo dựa vào là
$L_B$ giảm theo $B$, và mục~\ref{sec:mech-noise} kiểm điều đó bằng số.

\section{Cận co rút của gradient descent trên hàm bậc hai}

Với $H$ hằng và $t$ cố định, $w_{k+1} - \wstar = (I - tH)(w_k - \wstar)$. Khai
triển theo cơ sở riêng của $H$, thành phần ứng với trị riêng $\lambda_i$ nhân với
$1 - t\lambda_i$ mỗi vòng. Vì $f(w) - \fstar = \tfrac{1}{2}(w-\wstar)\trans H (w-\wstar)$,
sai số hàm mục tiêu nhân với $(1 - t\lambda_i)^{2}$ theo từng thành phần. Hệ số
tệ nhất là $\max_{\lambda_i \in [\mu, L]} \abs{1 - t\lambda_i}$, và giá trị $t$
làm nó nhỏ nhất là $t = 2/(L+\mu)$, khi ấy hai đầu khoảng cho cùng một giá trị
$(\kappa-1)/(\kappa+1)$. Bình phương lên được \eqref{eq:gd-rate}.

Chú ý rằng lập luận này cũng cho luôn điều kiện hội tụ: cần
$\abs{1 - t\lambda_i} < 1$ với mọi $i$, tức $0 < t < 2/L$, và tại $t = 2/L$ thì
thành phần ứng với $\lambda_i = L$ có hệ số bằng đúng $-1$ nên nó không tắt. Đó
là cơ chế mà mục~\ref{sec:mech-step} đo được.

\chapter{Cấu hình và seed}
\label{app:config}

Mọi con số thời gian trong báo cáo đo trên cùng một máy, cấu hình ghi ở
bảng~\ref{tab:config}, trong cùng một lượt chạy và không có tiến trình nặng nào
khác. Đổi máy thì các tỉ số giữa các phương
pháp cũng đổi, và mức đổi không đồng đều giữa các nhóm, nên không so trực tiếp
với bất kỳ con số nào đo ở nơi khác.

\begin{table}[tp]
  \centering
  \caption{Cấu hình phần cứng và phần mềm.}
  \label{tab:config}
  \begin{tabular}{ll}
    \toprule
    Hạng mục & Giá trị \\
    \midrule
    CPU        & Intel Core 7 150U, 10 nhân, 12 luồng \\
    Bộ nhớ     & 31 GB \\
    Hệ điều hành & Linux, glibc 2.39 \\
    Python     & 3.12.3 \\
    NumPy      & 2.5.2, liên kết OpenBLAS 0.3.34 đa luồng \\
    SciPy      & 1.18.0 \\
    scikit-learn & 1.9.0 \\
    pandas     & 3.0.5 \\
    matplotlib & 3.11.1 \\
    \bottomrule
  \end{tabular}
\end{table}

Mọi thành phần ngẫu nhiên dùng seed 0: phép lấy mẫu \num{200000} bản ghi từ một
triệu, phép chia tập huấn luyện và tập kiểm tra, phép chia fold của
cross-validation, và thứ tự lô của SGD. Bài toán đã chốt gồm $n = \num{160000}$
điểm huấn luyện, $d = 280$ thuộc tính, $\num{40000}$ điểm kiểm tra và
$\lambda = \num{0.01}$; ba hằng số $L$, $\mu$, $\kappa$ cùng $\fstar$ ghi ở
bảng~\ref{tab:constants}.

Mỗi cấu hình chạy ba lần và lấy trung vị theo thời gian. Chênh lệch giữa ba lần
chạy của cùng một cấu hình nằm trong khoảng 5\% tới 25\% ở phần lớn các nhóm và
lên tới 61\% ở trường hợp xấu nhất, nên chỉ những khoảng cách lớn hơn mức đó mới
nên đọc thành thứ hạng.

Kết quả từng nhóm ghi ra \texttt{results/raw/} dưới dạng JSON ngay khi nhóm ấy chạy
xong, nên vẽ lại được mọi hình mà không cần chạy lại thí nghiệm, và một lần gián
đoạn chỉ làm mất đúng nhóm đang chạy dở.
